\subsection{\bt{What executable reuse replaces}{可执行复用替代了哪些工作}}
\label{sec:sharing-eval}
\begin{figure}[t]
\centering
\ifbilingual
\includegraphics[width=\linewidth]{figures/reuse-work-zh}
\else
\includegraphics[width=\linewidth]{figures/reuse-work}
\fi
\caption{\bt{Agent work during reuse: mean tool calls over six sessions per method and database (three tasks, cutoffs 2--3), including failed attempts. Internal memory-model calls are excluded from this tool breakdown but included in Table~\ref{tab:official-results}'s token accounting.}{复用阶段的智能体工作：每种方法、每个数据库六个会话（三个任务、截点2--3）的平均工具调用数，保留失败尝试。此分解不计记忆模块内部的模型调用，其token已计入表~\ref{tab:official-results}。}}
\label{fig:reuse-work}
\Description{Two vertically stacked panels, one per database, compare all five methods. Stacked bars distinguish schema inspection, SQL probes, validation, memory tools, and submission with colors and hatching. The y-axis counts agent-issued tool calls per reuse session and starts at zero.}
\end{figure}

\bi{\emph{Agent work (Q2).} Figure~\ref{fig:reuse-work} decomposes the same RelBench reuse sessions into schema inspection, SQL probes, validation, memory tools, and submission. It includes failed attempts and counts calls initiated by the agent; the memories' internal model work is charged separately in Table~\ref{tab:official-results}. This distinguishes avoided exploration from retained checking. On rel-stack, \system\ makes \OfStackExplore\ schema/probe calls per reuse session versus \OfStackNoneExplore\ without memory, while still making \OfStackValidate\ validation calls. On rel-trial these counts are \OfTrialExplore\ versus \OfTrialNoneExplore\ exploration calls and \OfTrialValidate\ validations. Call counts describe the observed work distribution; they do not assign equal cost to every call or establish an isolated latency speedup.}
{\emph{智能体工作量（Q2）。}图~\ref{fig:reuse-work}将同一组RelBench复用会话分解为模式检查、SQL探查、验证、记忆工具与提交。统计保留失败尝试，计算智能体发起的调用；记忆模块内部的模型工作单独计入表~\ref{tab:official-results}。这样可以区分省去的探索与仍需执行的检查。rel-stack上，\system\ 每个复用会话进行\OfStackExplore{}次看结构或探查，无记忆为\OfStackNoneExplore{}次，同时仍进行\OfStackValidate{}次验证。rel-trial上分别为\OfTrialExplore{}与\OfTrialNoneExplore{}次探索，以及\OfTrialValidate{}次验证。这些调用数描述实际工作分布，不假设每次调用代价相等，也不构成隔离的延迟加速结论。}

\bi{\emph{Which knowledge was reused.} Each database records nine task-program lookups; rel-stack and rel-trial additionally record three and six feature-definition lookups. Neither sequence invokes the explicit feature-composition tool. Thus the end-to-end savings cannot be attributed specifically to feature composition: the evaluated policy combines task-program warm starts with optional definition retrieval. Publication logs provide a separate check on structured learning. Of 66 rel-stack drafts, 34 are newly promoted, nine corroborate existing definitions, and 23 fail replay; of nine rel-trial drafts, eight are promoted and one fails replay. Rejected drafts do not enter the shared feature library. A component ablation would be needed to separate full-program reuse from definition-level benefits.}
{\emph{实际复用了什么知识。}每个数据库记录了9次任务方案检索，rel-stack与rel-trial另分别记录3次和6次特征定义检索。两条序列均未调用显式特征组合工具。因此，不能将端到端节省单独归因于特征组合：受评估策略同时使用任务方案作为初始解，并按需检索定义。发布日志单独检验结构化学习：rel-stack的66个草稿中，34个新发布、9个为既有定义增加佐证、23个未通过重放；rel-trial的9个草稿中，8个发布、1个未通过重放。被拒绝的草稿不进入共享特征库。区分整份方案复用与定义级收益仍需组件消融。}

\bi{\emph{Publication fidelity in the metric studies.} The separate controlled metric workload also checks whether extraction preserves a verified answer. Across \PubAllJudged\ verified sessions, SQL replay rejects \PubReplayRefused\ erroneous candidates and the calculation check rejects \PubChainRefused\ more. In one example, extraction drops the electronics-category filter: a verified answer of \PubExAnswer\ becomes \PubExCanon\ under the candidate SQL, a \PubExRatio-fold error. This experiment tests the publication gate's role, independently of the RelBench token comparison.}
{\emph{指标实验中的发布保真性。}独立的受控指标负载检验提取是否保留已核实的答案。在\PubAllJudged{}个答案已核实的会话中，SQL重放拒绝\PubReplayRefused{}个错误候选，算式检查另拒绝\PubChainRefused{}个。一例中，提取丢失了电子品类过滤，使已验证答案\PubExAnswer{}变为候选SQL返回的\PubExCanon，误差达\PubExRatio{}倍。该实验检验发布检查的作用，与RelBench的token比较各自回答不同问题。}
